\input{sol_pre.tex}
\input{q_2024-12-01_2.tex}
\begin{document}
\input{sol_head.tex}

\Savol{1}
$8\cdot40^{45}\cdot36^2\cdot625^{25}\cdot216=2^3\cdot2^{135}5^{45}\cdot2^43^4\cdot5^{100}\cdot2^33^3=2^{145}5^{145}\cdot3^7=2187\cdot10^{145}$.\\
Raqamlar soni: $4+145=149$.
\Javob{A}

\Savol{2}
\[\frac{70+35+21+14+10}{2100}=\frac{150}{2100}=\frac1{14}.\]
\Javob{B}

\Savol{3}
$AB=S$. Qarama-qarşi harakat: $v_1+v_2=\dfrac S{t_1}$; bir tomonga (quvib yetiş): $v_1-v_2=\dfrac S{t_2}$.
\[v_1v_2=\frac{(v_1+v_2)^2-(v_1-v_2)^2}4=\frac{S^2}4\left(\frac1{t_1^2}-\frac1{t_2^2}\right).\]
\Javob{C}

\Savol{4}
$15\,000\,000\cdot0{,}4\cdot0{,}3=1\,800\,000$ söm.
\Javob{A}

\Savol{5}
$\text{EKUB}^2$ köpaytmaning böluvçisi: $\text{EKUB}\le3^2\cdot5\cdot7=315$ \ ($a=315$, $b=3^2\cdot5\cdot7\cdot11$).
\Javob{B}

\Savol{6}
$\tg\dfrac\pi4=1$, \ $\sin\dfrac{7\pi}3=\sin\dfrac\pi3=\dfrac{\sqrt3}2$: \ $1+\sqrt3\cdot\dfrac{\sqrt3}2=2{,}5$.
\Javob{D}

\Savol{7}
$3a_3=15$ $\Rightarrow$ $a_3=5$. \ $(5-d)^2+25+(5+d)^2=75+2d^2=83$ $\Rightarrow$ $d=-2$ (kamayuvçi).\\
$a_{13}=a_3+10d=5-20=-15$.
\Javob{C}

\Savol{8}
$54\pm30\sqrt3=\left(3\pm\sqrt3\right)^3$: \ $\left(3+\sqrt3\right)+\left(3-\sqrt3\right)=6$.
\Javob{C}

\Savol{9}
$x\ge4$: $x+11=9$ — yöq. \ $-2\le x<4$: $3x+3=9$, $x=2$. \ $-5\le x<-2$: $x-1=9$ — yöq. \ $x<-5$: $-x-11=9$, $x=-20$.\\
Ildizlar: $2$ va $-20$ — $2$ ta.
\Javob{B}

\Savol{10}
$b_4+b_5+b_6=q^3(b_1+b_2+b_3)$: \ $q^3=\dfrac{702}{26}=27$, $q=3$.
\Javob{D}

\Savol{11}
$\sqrt{25-x^2}$ ning qiymatlar töplami $[0;\,5]$.
\Javob{A}

\Savol{12}
$\ln7$ ni birlik deb olsak: $\ln6=a$, $\ln12=ab$ $\Rightarrow$ $\ln2=ab-a$, $\ln3=2a-ab$.
\[\log_{81}98\sqrt3=\frac{\ln2+2\ln7+\frac12\ln3}{4\ln3}=\frac{ab-a+2+a-\frac{ab}2}{4(2a-ab)}=\frac{4+ab}{16a-8ab}.\]
\Javob{A}

\Savol{13}
$t=x-\frac\pi3$: \ $(2\sin t-1)(\sin t-2)=0$ $\Rightarrow$ $\sin t=\frac12$, \ $t\in\left[-\frac{4\pi}3;\,\frac{5\pi}3\right]$.\\
$t=-\frac{7\pi}6,\ \frac\pi6,\ \frac{5\pi}6$ — $3$ ta.
\Javob{C}

\Savol{14}
$3^x-3^{x-2}=\frac89\cdot3^x$. \ $q=\left(\frac23\right)^x$ ($q\ne1$): \ $\dfrac{8/9}{1-q}>1+q$.\\
$q<1$: $\frac89>1-q^2$ $\Rightarrow$ $q>\frac13$; \ $q>1$ da çap tomon manfiy — yeçim yöq.\\
$\frac13<\left(\frac23\right)^x<1$ $\Rightarrow$ $x\in\left(0;\,\log_{\frac23}\frac13\right)$.
\Javob{B}

\Savol{15}
$x>1$. $u=\lg(x+1)$, $v=\lg(x-1)$: \ $u^2-uv-2v^2=(u-2v)(u+v)=0$.\\
$u=2v$: $x+1=(x-1)^2$ $\Rightarrow$ $x=3$. \ $u=-v$: $x^2-1=1$ $\Rightarrow$ $x=\sqrt2$. \ Köpaytma $3\sqrt2$.
\Javob{B}

\Savol{16}
$\dfrac1{5+2\sqrt6}=5-2\sqrt6$: \ $5-2\sqrt6-5+2\sqrt6=0$.
\Javob{C}

\Savol{17}
$t=\sqrt[4]x\ge0$: \ $t^2-5t+6\ge0$ $\Rightarrow$ $t\le2$ yoki $t\ge3$ $\Rightarrow$ $x\in[0;\,16]\cup[81;\,\infty)$.
\Javob{B}

\Savol{18}
\[\frac{4a^2-2ab+b^2}{2a-b}+b=\frac{4a^2}{2a-b};\qquad \frac{4a^2}{2a-b}\cdot\frac{2a-b}{a^2}=4.\]
\Javob{C}

\Savol{19}
Aniqlaniş sohasi: $x\le-1{,}5$. \ $\left(\sqrt{-3-2x}\right)^2=-3-2x$:
\[\frac{(2x-3)(x+2)}x\cdot\frac{-2(x+7)(x-2)}x\ge0.\]
$x\le-1{,}5$ da $(2x-3)(x-2)>0$ $\Rightarrow$ $(x+2)(x+7)\le0$, \ $-7\le x\le-2$.\\
Butun yeçimlar yiğindisi: $-7-6-5-4-3-2=-27$.
\Javob{C}

\Savol{20}
$t=\sqrt{9+x^2}\ge3$: \ $y=t^2-t-9$ — $t\ge3$ da ösuvçi, $y_{\min}=9-3-9=-3$.
\Javob{D}

\Savol{21}
$h-f=2$ — özgarmas $\Rightarrow$ $k=m$, $l-n=2$. \ $h\cdot f$ ning nollari $-3$ va $-1$, $h(0)f(0)=ln=3$:\\
$h=x+3$, $f=x+1$. \ $\dfrac{k+l}{m+n}=\dfrac42=2$.
\Javob{B}

\Savol{22}
$f'(x)=3^x\ln3\cdot\log_3(3x)+\dfrac{3^x}{x\ln3}$: \ $f'(1)=3\ln3+\dfrac3{\ln3}=\dfrac{3\ln^23+3}{\ln3}$.
\Javob{B}

\Savol{23}
Kesişiş: $x=-2$ va $x=0$ ($x+1=1$).
\[S=\int_{-2}^0(x+2)\,dx-\int_{-2}^0\left|x^2+3x+2\right|dx=2-\left(\frac16+\frac56\right)=1.\]
\Javob{B}

\Savol{24}
\fig{\begin{tikzpicture}[scale=0.42,font=\small]
\coordinate (O) at (0,0); \coordinate (A) at (-2,0); \coordinate (D) at (4,0); \coordinate (B) at (-2,6); \coordinate (C) at (4,6); \coordinate (K) at (-2,4); \coordinate (L) at (4,2);
\draw (A)--(B)--(C)--(D)--cycle; \draw (O) circle (4.472);
\draw[axs] (O)--(K); \draw[axs] (O)--(L); \fill (O) circle (2pt);
\node[below] at (O) {$O$}; \node[below left] at (A) {$A$}; \node[above left] at (B) {$B$}; \node[above right] at (C) {$C$}; \node[below right] at (D) {$D$};
\node[left] at (K) {$K$}; \node[right] at (L) {$L$}; \node[aux,below] at (-1,0) {$2$}; \node[aux,left] at (-2,2) {$4$};
\end{tikzpicture}}
$AO=6-4=2$. \ $R^2=OD^2+DL^2=20$ $\Rightarrow$ $AK=4$.\\
$\triangle OAK=\triangle LDO$ ($AO=DL=2$, $AK=OD=4$) $\Rightarrow$ $\angle AOK+\angle DOL=90^\circ$, \ $\angle KOL=90^\circ$.
\[\overset{\frown}{KL}=\frac14\cdot2\pi\cdot2\sqrt5=\sqrt5\,\pi.\]
\Javob{C}

\Savol{25}
$\angle C=90^\circ$ ($3^2+4^2=5^2$). \ $AH=\dfrac{AC^2}{AB}=\dfrac95$, \ $AL=AB\cdot\dfrac{AC}{AC+BC}=\dfrac{15}7$.
\[HL=\frac{15}7-\frac95=\frac{12}{35}.\]
\Javob{D}

\Savol{26}
\fig{\begin{tikzpicture}[scale=0.55,font=\small]
\coordinate (B) at (0,0); \coordinate (E) at (1,0); \coordinate (O) at (5,0); \coordinate (C) at (9,0); \coordinate (D) at (1.8,2.4); \coordinate (A) at (3.24,4.32);
\draw (A)--(B)--(C)--cycle; \draw (O) circle (4); \draw (D)--(C);
\draw[axs] (O)--(D); \coordinate (P) at ($(D)!0.5!(C)$); \draw[ax] (O)--(P);
\fill (O) circle (2pt);
\node[below left] at (B) {$B$}; \node[below] at (E) {$E$}; \node[below] at (O) {$O$}; \node[right] at (C) {$C$}; \node[above left] at (D) {$D$}; \node[above] at (A) {$A$}; \node[aux,above] at (P) {$P$};
\end{tikzpicture}}
$OD\perp AB$, $OD=OC=4$, $OB=5$ $\Rightarrow$ $BD=3$, \ $\sin B=\frac45$, $\cos B=\frac35$.\\
$BC=9$: \ $AC=\frac{36}5$, $AB=\frac{27}5$, $AD=\frac{12}5$. \ $DC=\sqrt{\frac{144+1296}{25}}=\frac{12\sqrt{10}}5$.\\
$\triangle ODC$ teng yonli: $OP\perp DC$, $DP=\frac{6\sqrt{10}}5$.
\[OP=\sqrt{16-\frac{360}{25}}=\frac{2\sqrt{10}}5.\]
\Javob{A}

\Savol{27}
\fig{\begin{tikzpicture}[scale=0.42,font=\small]
\coordinate (A) at (0,0); \coordinate (B) at (0,4); \coordinate (C) at (5,4); \coordinate (D) at (8,0); \coordinate (H) at (5,0);
\draw (A)--(B)--(C)--(D)--cycle; \draw (B)--(D); \draw[ax] (C)--(H);
\node[below left] at (A) {$A$}; \node[above left] at (B) {$B$}; \node[above] at (C) {$C$}; \node[below right] at (D) {$D$}; \node[aux,below] at (H) {$H$};
\node[left] at (0,2) {$4$}; \node[aux,above] at (2.5,4) {$a$}; \node[aux,right] at (6.5,2) {$a$};
\end{tikzpicture}}
$AD=\sqrt{80-16}=8$. \ $BC=CD=a$, $HD=8-a$: \ $16+(8-a)^2=a^2$ $\Rightarrow$ $a=5$.
\[S=\frac{5+8}2\cdot4=26.\]
\Javob{C}

\Savol{28}
\fig{\begin{tikzpicture}[scale=0.5,font=\small]
\coordinate (H) at (0,0); \coordinate (A) at (0,6); \coordinate (M) at (8,0); \coordinate (P) at (2.88,3.84);
\draw (H)--(M)--(A)--cycle; \draw[axs] (H)--(P);
\draw ($(H)!0.5cm!(M)$)--($(H)!0.5cm!(M)!0.5cm!90:(M)$)--($(H)!0.5cm!(A)$);
\node[below left] at (H) {$H$}; \node[above] at (A) {$A$}; \node[below right] at (M) {$M$}; \node[aux,above right] at (P) {$P$};
\node[left] at (0,3) {$6$}; \node at (6.6,0.4) {\scriptsize$37^\circ$};
\end{tikzpicture}}
$\angle HAM=53^\circ$: \ $HP=AH\sin53^\circ=AH\cos37^\circ=6\cdot0{,}8=4{,}8$.
\Javob{C}

\Savol{29}
$E$ nuqta $AB$ ning $B$ dan taşqari davomida, $BE=\frac12AB$ $\Rightarrow$ $AE=\frac32AB$.\\
$C$ dan balandlik umumiy: $S_{ACE}=\frac32\cdot60=90$.
\Javob{C}

\Savol{30}
Uçdan çiquvçi qöşni diagonallar orasidagi burçak $\frac{180^\circ}n$ (içki çizilgan burçak). Eng kiçik diagonal $2$-uçga, eng kattasi (diametr) $\frac n2$-uçga:
\[\left(\frac n2-2\right)\frac{180^\circ}n=45^\circ\ \Rightarrow\ n=8;\qquad 180^\circ\cdot6=1080^\circ.\]
\Javob{A}

\Savol{31}
1) $x^2=\sqrt{10}-\pi>0$ — böş emas. \ 2), 3) — böş. \ 4) $x=-1,\,-2$ — irratsional emas, böş. \ Jami $3$ ta.
\Javob{C}

\Savol{32}
Hammasi $9000$ ta. Oxirgi raqam $5$, qolgan uçtasi toq: $5^3=125$. \ $P=\dfrac{125}{9000}=\dfrac1{72}$.
\Javob{A}

\newpage
\Savol{33}
\fig{\begin{tikzpicture}[scale=0.8,font=\small]
\coordinate (O) at (0,0); \coordinate (A) at (0,2.1); \coordinate (D) at (1.6,-1.6); \coordinate (C) at (2.6,0.7); \coordinate (B) at (2.1,-0.45);
\draw[-{Stealth[length=4pt]}] (O)--(0,3) node[left] {$z$}; \draw[-{Stealth[length=4pt]}] (O)--(2.1,-2.1) node[right] {$x$}; \draw[-{Stealth[length=4pt]}] (O)--(3.6,0.97) node[right] {$y$};
\fill[okgreen!30] (O)--(D)--(C)--cycle; \draw (D)--(C); \draw[axs] (O)--(B); \draw (A)--(D); \draw (A)--(C); \draw[dashed] (A)--(B);
\fill (A) circle (1.5pt); \fill (B) circle (1.5pt);
\node[left] at (A) {$A$}; \node[left] at (O) {$O$}; \node[below left] at (D) {$D$}; \node[above] at (C) {$C$}; \node[right] at (B) {$B$};
\end{tikzpicture}}
$OB=\sqrt{2+2}=2$. $OB$ balandlik va mediana $\Rightarrow$ $OC=OD$, $B$ — $CD$ ning örtasi: $D\left(2\sqrt2;0;0\right)$, $C\left(0;2\sqrt2;0\right)$, $BC=BD=2$.

\textbf{33.} $V=\dfrac13\cdot\dfrac12\cdot2\sqrt2\cdot2\sqrt2\cdot3=4$ — \textbf{B}.

\textbf{34.} $AO$ atrofida — konus: $r=OB=2$, $h=AO=3$: \ $V=\dfrac13\pi\cdot4\cdot3=4\pi=12$ — \textbf{D}.

\textbf{35.} $CD$ atrofida — ikkita konus: $r=OB=2$, $h=BC=BD=2$:\\
$V=2\cdot\dfrac13\pi\cdot4\cdot2=\dfrac{16\pi}3=16$ — \textbf{E}.
\Javob{33 — B, \ 34 — D, \ 35 — E}

\newpage
\Savol{36}
Viyet: $a+b+c=9$, $ab+bc+ca=11$, $abc=1$ (ildizlar musbat).

a) $9+1+11=21$.

b) $s=\sqrt{ab}+\sqrt{bc}+\sqrt{ca}$:
\[t^2=a+b+c+2s=9+2s,\qquad s^2=ab+bc+ca+2\sqrt{abc}\,t=11+2t.\]
$\left(t^2-9\right)^2=4s^2=44+8t$ $\Rightarrow$ $t^4-18t^2+81=44+8t$:
\[t^4-18t^2-8t=-37.\]
\Javob{a) 21; \quad b) $-37$}

\newpage
\Savol{37}
$u=\cos^2\frac{\pi x}{18}$: \ $\sqrt{7+8u}=4(1-u)$ $\Rightarrow$ $7+8u=16(1-u)^2$:
\[16u^2-40u+9=0\ \Rightarrow\ u=\frac14\ \ \left(u=\frac94>1\right).\]
$\cos^2\frac{\pi x}{18}=\frac14$ $\Rightarrow$ $\cos\frac{\pi x}9=-\frac12$ $\Rightarrow$ $\frac{\pi x}9=\pm\frac{2\pi}3+2\pi n$, \ $x=\pm6+18n$.

a) $[0;\,18]$: $x=6,\ 12$ — yiğindi $18$.

b) $[20;\,30]$: $x=24,\ 30$ — $2$ ta.
\Javob{a) 18; \quad b) 2 ta}

\newpage
\Savol{38}
$y=\dfrac{x+3}{x-1}$ $\Rightarrow$ $xy-y=x+3$ $\Rightarrow$ $x=\dfrac{y+3}{y-1}$: \ $f^{-1}(x)=\dfrac{x+3}{x-1}=f(x)$.

a) $\dfrac{x+3}{x-1}=2$ $\Rightarrow$ $x=5$.

b) $f(x)=f^{-1}(x)$ aniqlaniş sohasidagi barça $x$ da örinli; ildiz böla olmaydigan butun son faqat $x=1$ — $1$ ta.
\Javob{a) 5; \quad b) 1 ta}

\newpage
\Savol{39}
\fig{\begin{tikzpicture}[scale=0.25,font=\small]
\draw (0,0)--(16,0)--(6,15)--(0,15)--cycle; \fill[okgreen!30] (0,0) rectangle (8,12); \draw (0,0) rectangle (8,12);
\draw (0,0)--(16,0)--(6,15)--(0,15)--cycle;
\node[below] at (8,0) {$16$}; \node[above] at (3,15) {$6$}; \node[left] at (0,7.5) {$15$};
\node[aux,right] at (8,6) {$h$}; \node[aux,above] at (4,12) {$w$};
\end{tikzpicture}}
Töğri törtburçakning balandligi $h$. $h$ balandlikda trapetsiya kengligi $16-\frac{10}{15}h=16-\frac23h$ (yon tomon çiziği böyiça).
\[S(h)=h\left(16-\frac23h\right),\qquad S'(h)=16-\frac43h=0\ \Rightarrow\ h=12,\ w=8.\]

a) $S=12\cdot8=96$.

b) $S_{\text{trap}}=\dfrac{6+16}2\cdot15=165$: \ $165-96=69$.
\Javob{a) 96; \quad b) 69}

\newpage
\Savol{40}
$f'(x)=4x^3+3x^2+C$, \ $f'(1)=7+C=9$ $\Rightarrow$ $C=2$.\\
$f(x)=x^4+x^3+2x+D$, \ $f(1)=4+D=5$ $\Rightarrow$ $D=1$.

a) $f(0)=1$.

b) $g(x)+x\,g'(x)=\left(x\,g(x)\right)'$:
\[\int_{2023}^{2024}\left(x\,g(x)\right)'dx=2024\cdot2-2023\cdot3=-2021.\]
\Javob{a) 1; \quad b) $-2021$}

\newpage
\Savol{41}
\fig{\begin{tikzpicture}[scale=0.85,font=\small]
\coordinate (B) at (0,0); \coordinate (C) at (2,0); \coordinate (A) at (0,-2); \coordinate (D) at (2,-2); \coordinate (E) at (-1.414,1.414);
\fill[okgreen!30] (E)--(C) arc (0:-90:2) --cycle;
\draw (B) circle (2); \draw (A)--(B)--(C)--(D)--cycle; \draw (E)--(C); \draw (E)--(A); \draw[ax] (E)--(D);
\node[above left] at (E) {$E$}; \node[above left] at (B) {$B$}; \node[right] at (C) {$C$}; \node[below] at (A) {$A$}; \node[right] at (D) {$D$};
\draw (0:0.35) arc (0:135:0.35); \node at (0.55,0.5) {\scriptsize$135^\circ$};
\end{tikzpicture}}
Kvadrat tomoni $r$. Kvadrat içida böyalgani — çorak doira, böyalmagani — qolgan qism.

a) $\dfrac{\pi r^2/4}{r^2-\pi r^2/4}=\dfrac\pi{4-\pi}$.

b) $\angle EBC=\angle EBA=135^\circ$. Böyalgan soha $=$ çorak doira $+$ $\triangle EBC$ $+$ $\triangle EBA$:
\[S=\frac{\pi\cdot4}4+2\cdot\frac12\cdot2\cdot2\cdot\sin135^\circ=\pi+2\sqrt2.\]
\Javob{a) $\dfrac\pi{4-\pi}$; \quad b) $2\sqrt2+\pi$}

\newpage
\Savol{42}
\fig{\begin{tikzpicture}[scale=0.32,font=\small]
\coordinate (C) at (0,0); \coordinate (A) at (0,14); \coordinate (B) at (10,0); \coordinate (D) at (0,4); \coordinate (E) at (6,0); \coordinate (O) at (5.172,1.931);
\draw (A)--(B)--(C)--cycle; \draw (A)--(E); \draw (D)--(B); \draw (0,0.8)--(0.8,0.8)--(0.8,0);
\node[above] at (A) {$A$}; \node[below right] at (B) {$B$}; \node[below left] at (C) {$C$}; \node[left] at (D) {$D$}; \node[below] at (E) {$E$}; \node[above right] at (O) {$O$};
\node[aux,left] at (0,9) {$a$}; \node[aux,left] at (0,2) {$d$}; \node[aux,below] at (8,0) {$d$};
\node at (0.45,11.6) {\scriptsize$\alpha$}; \node at (8.4,0.5) {\scriptsize$\beta$};
\end{tikzpicture}}
$BC=AD=a$, $CD=BE=d$: \ $AC=a+d$, $CE=a-d$. \ $\alpha=\angle CAE$, $\beta=\angle CBD$:
\[\tg\alpha=\frac{a-d}{a+d},\quad \tg\beta=\frac da,\quad \tg(\alpha+\beta)=\frac{a(a-d)+d(a+d)}{a(a+d)-d(a-d)}=\frac{a^2+d^2}{a^2+d^2}=1.\]
$\alpha+\beta=45^\circ$.

a) $\triangle BCD$ da $\angle CDB=90^\circ-\beta$ $\Rightarrow$ $\angle ADO=90^\circ+\beta$. \ $\triangle AOD$ da:
\[\angle AOD=180^\circ-\alpha-(90^\circ+\beta)=90^\circ-45^\circ=45^\circ.\]

b) $a=10$: \ $AE^2=(10+d)^2+(10-d)^2=200+2d^2=232$ $\Rightarrow$ $d=4$, $AC=14$.
\[S_{ABC}=\frac12\cdot14\cdot10=70.\]
\Javob{a) $45^\circ$; \quad b) 70}

\newpage
\Savol{43}
\fig{\begin{tikzpicture}[scale=0.75,font=\small]
\coordinate (A) at (0,0); \coordinate (D) at (6,0); \coordinate (B) at (2.8,3); \coordinate (C) at (8.8,3);
\coordinate (A1) at (1.76,0.6); \coordinate (C1) at (7.04,2.4); \coordinate (E) at (5.36,0.6); \coordinate (B1) at (3.12,2.7); \coordinate (D1) at (5.253,0.7);
\draw (A)--(B)--(C)--(D)--cycle; \draw (A)--(C); \draw (B)--(D);
\fill[okgreen!30] (A1)--(C1)--(E)--cycle; \draw (A1)--(C1)--(E)--cycle;
\draw[axs] (A1)--(B1)--(C1)--(D1)--cycle;
\node[below left] at (A) {$A$}; \node[below] at (D) {$D$}; \node[above left] at (B) {$B$}; \node[above right] at (C) {$C$};
\node[below] at (A1) {$A_1$}; \node[above] at (C1) {$C_1$}; \node[below right] at (E) {$E$}; \node[above] at (B1) {$B_1$}; \node[left] at (D1) {$D_1$};
\end{tikzpicture}}
a) $A_1E\parallel AD$, $C_1E\parallel CD$, $A_1C_1$ esa $AC$ ustida $\Rightarrow$ $\triangle A_1C_1E\sim\triangle ACD$, \ $k=\dfrac{A_1C_1}{AC}=\dfrac35$:
\[\frac{S_{A_1C_1E}}{S_{ACD}}=\left(\frac35\right)^2=\frac9{25}.\]

b) Diagonallari $d_1$, $d_2$ va ular orasidagi burçagi $\varphi$ bölgan törtburçak yuzi $\frac12d_1d_2\sin\varphi$. $A_1B_1C_1D_1$ ning diagonallari $AC$ va $BD$ ustida yotadi ($\varphi$ umumiy):
\[\frac{S_{A_1B_1C_1D_1}}{S_{ABCD}}=\frac{A_1C_1}{AC}\cdot\frac{B_1D_1}{BD}=\frac35\cdot\frac23=\frac25.\]
\Javob{a) $\dfrac9{25}$; \quad b) $\dfrac25$}

\newpage
\Savol{44}
\fig{\begin{tikzpicture}[scale=1,font=\small]
\draw (0,0) circle (1.5); \draw[dashed] (0,0)--(1.5,0);
\coordinate (S) at (0,1.5); \coordinate (P) at (0,-0.9); \coordinate (Q) at (1.2,-0.9);
\draw (-1.2,-0.9)--(S)--(Q); \draw (0,-0.9) ellipse (1.2 and 0.28);
\draw[axs] (S)--(P)--(Q); \draw[ax] (0,0)--(Q); \fill (0,0) circle (1.3pt);
\node[above] at (S) {$S$}; \node[aux,left] at (0,0.75) {$H$}; \node[aux,below] at (0.6,-0.95) {$r$}; \node[left] at (0,0) {$O$}; \node[aux,above] at (0.75,0) {$R$};
\end{tikzpicture}}
$r^2=H(2R-H)$ (içki çizilgan konus). Şart:
\[\frac13\pi r^2H=\frac14\cdot\frac43\pi R^3\ \Rightarrow\ H^2(2R-H)=R^3.\]
$k=\dfrac HR$: \ $k^3-2k^2+1=0$ $\Rightarrow$ $(k-1)\left(k^2-k-1\right)=0$, \ $k=1$ yoki $k=\dfrac{1+\sqrt5}2$.

a) $H>R$ $\Rightarrow$ $k=\frac{1+\sqrt5}2$, \ $k^3=2k+1=2+\sqrt5$. \ $R^3=\dfrac{H^3}{k^3}=\dfrac{2+\sqrt5}{2+\sqrt5}=1$:
\[V_{\text{şar}}=\frac43\pi R^3=4.\]

b) $H=R=\sqrt{3+2\sqrt2}=1+\sqrt2$ $\Rightarrow$ $r^2=H\cdot R=H^2$, $r=H$, yasovçi $l=H\sqrt2$.\\
Öq kesimi — teng yonli uçburçak; içki aylana radiusi:
\[\rho=\frac{S}{p}=\frac{rH}{r+l}=\frac{H^2}{H\left(1+\sqrt2\right)}=\frac{1+\sqrt2}{1+\sqrt2}=1,\qquad S_{\text{sfera}}=4\pi\rho^2=12.\]
\Javob{a) 4; \quad b) 12}

\newpage
\Savol{45}
\fig{\begin{tikzpicture}[scale=0.38,font=\small]
\draw (-4,0)--(21,0); \draw[thick,domain=0:20,samples=60] plot (\x,{8-0.08*(\x-10)^2});
\draw[axs] (-2.5,0)--(5,6); \draw[ax] (5,0)--(5,6); \draw[dashed] (10,0)--(10,8);
\draw (-2.5,0)+(0:1.2) arc (0:38.66:1.2); \node at (-0.7,0.45) {\scriptsize$\alpha$};
\node[below] at (-2.5,0) {$D$}; \node[below] at (0,0) {$O$}; \node[below] at (5,0) {$B$}; \node[below] at (10,0) {$E$}; \node[below] at (20,0) {$K$};
\node[above left] at (5,6) {$A$}; \node[above] at (10,8) {$C$}; \node[right] at (5,3) {$6$}; \node[right] at (10,4) {$8$};
\end{tikzpicture}}
Karkas — $C$ uçli parabola; $E$ dan $t$ masofadagi balandlik:
\[y=8-a t^2,\quad t=10\ \text{da}\ y=0\ \Rightarrow\ a=0{,}08.\]
$A$: $8-0{,}08t^2=6$ $\Rightarrow$ $t=5$, ya'ni $BE=5$, $OB=5$.

a) Urinma oğmasi: $y'=-0{,}16t$, $t=-5$ da ($A$ çap tomonda) $0{,}16\cdot5=0{,}8$. \ $\tg\alpha=0{,}8$.

b) $\triangle ABD$: $DB=\dfrac{AB}{\tg\alpha}=\dfrac6{0{,}8}=7{,}5$. \ $OD=DB-OB=7{,}5-5=2{,}5$.
\Javob{a) 0,8; \quad b) 2,5}

\ShowJavoblar
\end{document}
